%% Prague Cafe Project -- Batch 1 slides (Beamer). Numbers come from ../report/tables/numbers.tex,
%% which analysis/03_tests.R writes. Graphs are read from figures/ (A1_figure.png ... A10_figure.png),
%% which analysis/03_tests.R writes into this folder. Compile from this folder:  latexmk -pdf slides.tex
\documentclass[aspectratio=169,11pt]{beamer}
\usepackage[T1]{fontenc}
\usepackage{mathpazo}
\usepackage{graphicx}
\usepackage{booktabs}

\definecolor{espresso}{HTML}{2B1D16}
\definecolor{copper}{HTML}{B8622B}
\definecolor{pale}{HTML}{E7E2DC}
\definecolor{greyc}{HTML}{8A7F76}

% macros that numbers.tex expects
\newcommand{\pval}[1]{$p$\,=\,#1}
\newcommand{\plt}[1]{$p$\,<\,#1}
\newcommand{\CI}[2]{[#1,\,#2]}
% numbers.tex -- written by analysis/03_tests.R; do not edit by hand
\newcommand{\Nall}{200}
\newcommand{\AoneBench}{90}
\newcommand{\AoneN}{200}
\newcommand{\AoneMean}{82.8}
\newcommand{\AoneLo}{80.7}
\newcommand{\AoneHi}{84.8}
\newcommand{\AoneT}{\ensuremath{-}6.94}
\newcommand{\AoneDf}{199}
\newcommand{\AoneP}{\plt{.001}}
\newcommand{\AoneD}{\ensuremath{-}0.49}
\newcommand{\AoneDLo}{\ensuremath{-}0.64}
\newcommand{\AoneDHi}{\ensuremath{-}0.34}
\newcommand{\AoneWmed}{82.5}
\newcommand{\AoneWP}{\plt{.001}}
\newcommand{\AoneBLo}{80.8}
\newcommand{\AoneBHi}{84.8}
\newcommand{\AtwoNa}{67}
\newcommand{\AtwoNb}{133}
\newcommand{\AtwoMa}{86.6}
\newcommand{\AtwoMb}{80.8}
\newcommand{\AtwoDiff}{5.7}
\newcommand{\AtwoLo}{0.9}
\newcommand{\AtwoHi}{10.5}
\newcommand{\AtwoD}{0.39}
\newcommand{\AtwoDLo}{0.10}
\newcommand{\AtwoDHi}{0.69}
\newcommand{\AtwoT}{2.37}
\newcommand{\AtwoDf}{102.5}
\newcommand{\AtwoP}{\pval{.010}}
\newcommand{\AtwoWP}{\pval{.069}}
\newcommand{\AtwoHL}{3.0}
\newcommand{\AtwoHLLo}{0.0}
\newcommand{\AtwoHLHi}{8.0}
\newcommand{\AtwoRB}{0.13}
\newcommand{\AtwoRBLo}{\ensuremath{-}0.05}
\newcommand{\AtwoRBHi}{0.30}
\newcommand{\AthreeN}{200}
\newcommand{\AthreeM}{18.1}
\newcommand{\AthreeLo}{16.9}
\newcommand{\AthreeHi}{19.2}
\newcommand{\AthreeDz}{2.19}
\newcommand{\AthreeDzLo}{1.93}
\newcommand{\AthreeDzHi}{2.44}
\newcommand{\AthreeT}{30.96}
\newcommand{\AthreeDf}{199}
\newcommand{\AthreeP}{\plt{.001}}
\newcommand{\AthreeWP}{\plt{.001}}
\newcommand{\AthreeHL}{17.5}
\newcommand{\AthreeHLLo}{17.5}
\newcommand{\AthreeHLHi}{18.5}
\newcommand{\AthreeSD}{8.3}
\newcommand{\AfourN}{190}
\newcommand{\AfourF}{2.33}
\newcommand{\AfourDfa}{5}
\newcommand{\AfourDfb}{58.8}
\newcommand{\AfourP}{\pval{.054}}
\newcommand{\AfourEta}{0.062}
\newcommand{\AfourEtaLo}{0.000}
\newcommand{\AfourEtaHi}{0.116}
\newcommand{\AfourOmega}{0.036}
\newcommand{\AfourH}{9.51}
\newcommand{\AfourKWdf}{5}
\newcommand{\AfourKWP}{\pval{.090}}
\newcommand{\AfourEps}{0.050}
\newcommand{\AfourBestPair}{Prague 1 minus Prague 7}
\newcommand{\AfourBestDiff}{8.5}
\newcommand{\AfourBestLo}{\ensuremath{-}0.2}
\newcommand{\AfourBestHi}{17.3}
\newcommand{\AfourBestP}{\pval{.062}}
\newcommand{\AfourNsigGH}{0}
\newcommand{\AfourNsigDunn}{0}
\newcommand{\AfourMax}{Prague 1}
\newcommand{\AfourMin}{Prague 3}
\newcommand{\AfourMaxM}{86.6}
\newcommand{\AfourMinM}{74.9}
\newcommand{\AfiveN}{199}
\newcommand{\AfiveChi}{5.31}
\newcommand{\AfiveDf}{2}
\newcommand{\AfiveP}{\pval{.070}}
\newcommand{\AfiveFisherP}{\pval{.074}}
\newcommand{\AfiveMinExp}{3.70}
\newcommand{\AfiveV}{0.16}
\newcommand{\AfiveVLo}{0.05}
\newcommand{\AfiveVHi}{0.31}
\newcommand{\AfiveShareInd}{33.6}
\newcommand{\AfiveShareLoc}{27.5}
\newcommand{\AfiveShareInt}{63.6}
\newcommand{\AfiveNInt}{11}
\newcommand{\AfiveKInt}{7}
\newcommand{\AfiveResInt}{2.16}
\newcommand{\AfiveResInd}{\ensuremath{-}0.04}
\newcommand{\AfiveResLoc}{\ensuremath{-}1.09}
\newcommand{\AsixN}{200}
\newcommand{\AsixR}{\ensuremath{-}0.17}
\newcommand{\AsixRLo}{\ensuremath{-}0.30}
\newcommand{\AsixRHi}{\ensuremath{-}0.03}
\newcommand{\AsixRsq}{2.7}
\newcommand{\AsixT}{\ensuremath{-}2.36}
\newcommand{\AsixDf}{198}
\newcommand{\AsixP}{\pval{.010}}
\newcommand{\AsixRho}{\ensuremath{-}0.18}
\newcommand{\AsixRhoLo}{\ensuremath{-}0.31}
\newcommand{\AsixRhoHi}{\ensuremath{-}0.03}
\newcommand{\AsixRhoP}{\pval{.006}}
\newcommand{\AsixTau}{\ensuremath{-}0.12}
\newcommand{\AsixTauLo}{\ensuremath{-}0.23}
\newcommand{\AsixTauHi}{\ensuremath{-}0.02}
\newcommand{\AsixTauP}{\pval{.006}}
\newcommand{\AsevenN}{200}
\newcommand{\AsevenB}{\ensuremath{-}1.85}
\newcommand{\AsevenLo}{\ensuremath{-}3.74}
\newcommand{\AsevenHi}{0.05}
\newcommand{\AsevenP}{\pval{.057}}
\newcommand{\AsevenT}{\ensuremath{-}1.92}
\newcommand{\AsevenDf}{198}
\newcommand{\AsevenRsq}{0.027}
\newcommand{\AsevenLogB}{\ensuremath{-}0.020}
\newcommand{\AsevenPct}{\ensuremath{-}2.0}
\newcommand{\AsevenPctLo}{\ensuremath{-}4.2}
\newcommand{\AsevenPctHi}{0.1}
\newcommand{\AsevenLogP}{\pval{.067}}
\newcommand{\AsevenLogT}{\ensuremath{-}1.84}
\newcommand{\AsevenLogRsq}{0.023}
\newcommand{\AsevenNinfl}{11}
\newcommand{\AsevenBinfl}{\ensuremath{-}1.21}
\newcommand{\AsevenMaxCook}{0.39}
\newcommand{\AsevenMaxDist}{8.1}
\newcommand{\AeightN}{199}
\newcommand{\AeightBa}{\ensuremath{-}1.84}
\newcommand{\AeightBb}{\ensuremath{-}1.77}
\newcommand{\AeightBbLo}{\ensuremath{-}3.58}
\newcommand{\AeightBbHi}{0.03}
\newcommand{\AeightBaLo}{\ensuremath{-}3.74}
\newcommand{\AeightBaHi}{0.06}
\newcommand{\AeightChange}{0.07}
\newcommand{\AeightRelChange}{3.5}
\newcommand{\AeightF}{6.67}
\newcommand{\AeightDfa}{2}
\newcommand{\AeightDfb}{195}
\newcommand{\AeightP}{\pval{.002}}
\newcommand{\AeightDeltaR}{0.062}
\newcommand{\AeightRsqA}{0.027}
\newcommand{\AeightRsqB}{0.090}
\newcommand{\AeightVifDist}{1.02}
\newcommand{\AeightVifType}{1.02}
\newcommand{\AeightGVIFadj}{1.01}
\newcommand{\AeightLocal}{6.1}
\newcommand{\AeightLocalLo}{1.9}
\newcommand{\AeightLocalHi}{10.3}
\newcommand{\AeightInt}{12.8}
\newcommand{\AeightIntLo}{\ensuremath{-}4.7}
\newcommand{\AeightIntHi}{30.4}
\newcommand{\AeightNinfl}{13}
\newcommand{\AeightBinfl}{\ensuremath{-}1.47}
\newcommand{\AnineN}{199}
\newcommand{\AnineNind}{137}
\newcommand{\AnineNch}{62}
\newcommand{\AnineInt}{\ensuremath{-}1.56}
\newcommand{\AnineIntLo}{\ensuremath{-}6.83}
\newcommand{\AnineIntHi}{3.70}
\newcommand{\AnineF}{1.00}
\newcommand{\AnineDfa}{1}
\newcommand{\AnineDfb}{195}
\newcommand{\AnineP}{\pval{.318}}
\newcommand{\AnineDeltaR}{0.005}
\newcommand{\AnineSlopeInd}{\ensuremath{-}1.33}
\newcommand{\AnineSlopeIndLo}{\ensuremath{-}2.95}
\newcommand{\AnineSlopeIndHi}{0.30}
\newcommand{\AnineSlopeCh}{\ensuremath{-}2.89}
\newcommand{\AnineSlopeChLo}{\ensuremath{-}7.90}
\newcommand{\AnineSlopeChHi}{2.12}
\newcommand{\AnineMaxDiff}{6.83}
\newcommand{\AtenN}{199}
\newcommand{\AtenEvents}{98}
\newcommand{\AtenAME}{\ensuremath{-}4.9}
\newcommand{\AtenAMELo}{\ensuremath{-}10.0}
\newcommand{\AtenAMEHi}{0.2}
\newcommand{\AtenOR}{0.81}
\newcommand{\AtenORLo}{0.65}
\newcommand{\AtenORHi}{1.02}
\newcommand{\AtenZ}{\ensuremath{-}1.81}
\newcommand{\AtenP}{\pval{.070}}
\newcommand{\AtenMcF}{0.039}
\newcommand{\AtenSampleMedian}{82.5}
\newcommand{\AtenShareHigh}{49.2}
\newcommand{\RobMinR}{\ensuremath{-}0.19}
\newcommand{\RobMaxR}{\ensuremath{-}0.12}
\newcommand{\ClnEmpty}{99}
\newcommand{\ClnDistrict}{20}
\newcommand{\ClnType}{1}
\newcommand{\ClnCap}{12}
\newcommand{\ClnEsp}{12}
\newcommand{\ClnSep}{6}
\newcommand{\ClnDup}{1}
\newcommand{\Nraw}{300}
\newcommand{\Nfilled}{201}
\newcommand{\Ntype}{199}
\newcommand{\PriceMean}{82.8}
\newcommand{\PriceSD}{14.8}
\newcommand{\PriceMedian}{82.5}
\newcommand{\PriceIQR}{15.0}
\newcommand{\PriceSkew}{0.53}
\newcommand{\PriceMin}{40}
\newcommand{\PriceMax}{140}
\newcommand{\DistSkew}{1.14}
\newcommand{\DistMedian}{1.64}
\newcommand{\DistMax}{8.1}
\newcommand{\DiffSkew}{0.64}
\newcommand{\AsevenSkewLin}{0.45}
\newcommand{\AsevenSkewLog}{\ensuremath{-}0.34}
\newcommand{\AoneDiffAbs}{7.2}
\newcommand{\AoneDir}{below}
\newcommand{\AonePctRel}{8.0}
\newcommand{\NstudentsCafes}{20}
\newcommand{\AthreePctEsp}{28}
\newcommand{\EspMean}{64.7}
\newcommand{\Ntests}{56}


\usetheme{default}
\usefonttheme{serif}
\setbeamertemplate{navigation symbols}{}
\setbeamercolor{background canvas}{bg=white}
\setbeamercolor{title}{fg=espresso}
\setbeamercolor{frametitle}{fg=espresso}
\setbeamercolor{normal text}{fg=espresso}
\setbeamercolor{itemize item}{fg=copper}
\setbeamerfont{frametitle}{series=\bfseries,size=\Large}
\setbeamertemplate{frametitle}{\vspace{0.4em}\insertframetitle}
\setbeamertemplate{footline}{\hfill\color{greyc}\small\insertframenumber\hspace{1em}\vspace{0.6em}}
\setbeamertemplate{itemize items}[circle]

\graphicspath{{figures/}}   % all graphs live in presentation/figures/  (A1_figure.png ... A10_figure.png)

% how the notes (the spoken script) are shown:
%   pdflatex slides.tex                                         -> slides only
%   pdflatex -jobname=slides_script "\def\withscript{1}%% Prague Cafe Project -- Batch 1 slides (Beamer). Numbers come from ../report/tables/numbers.tex,
%% which analysis/03_tests.R writes. Graphs are read from figures/ (A1_figure.png ... A10_figure.png),
%% which analysis/03_tests.R writes into this folder. Compile from this folder:  latexmk -pdf slides.tex
\documentclass[aspectratio=169,11pt]{beamer}
\usepackage[T1]{fontenc}
\usepackage{mathpazo}
\usepackage{graphicx}
\usepackage{booktabs}

\definecolor{espresso}{HTML}{2B1D16}
\definecolor{copper}{HTML}{B8622B}
\definecolor{pale}{HTML}{E7E2DC}
\definecolor{greyc}{HTML}{8A7F76}

% macros that numbers.tex expects
\newcommand{\pval}[1]{$p$\,=\,#1}
\newcommand{\plt}[1]{$p$\,<\,#1}
\newcommand{\CI}[2]{[#1,\,#2]}
% numbers.tex -- written by analysis/03_tests.R; do not edit by hand
\newcommand{\Nall}{200}
\newcommand{\AoneBench}{90}
\newcommand{\AoneN}{200}
\newcommand{\AoneMean}{82.8}
\newcommand{\AoneLo}{80.7}
\newcommand{\AoneHi}{84.8}
\newcommand{\AoneT}{\ensuremath{-}6.94}
\newcommand{\AoneDf}{199}
\newcommand{\AoneP}{\plt{.001}}
\newcommand{\AoneD}{\ensuremath{-}0.49}
\newcommand{\AoneDLo}{\ensuremath{-}0.64}
\newcommand{\AoneDHi}{\ensuremath{-}0.34}
\newcommand{\AoneWmed}{82.5}
\newcommand{\AoneWP}{\plt{.001}}
\newcommand{\AoneBLo}{80.8}
\newcommand{\AoneBHi}{84.8}
\newcommand{\AtwoNa}{67}
\newcommand{\AtwoNb}{133}
\newcommand{\AtwoMa}{86.6}
\newcommand{\AtwoMb}{80.8}
\newcommand{\AtwoDiff}{5.7}
\newcommand{\AtwoLo}{0.9}
\newcommand{\AtwoHi}{10.5}
\newcommand{\AtwoD}{0.39}
\newcommand{\AtwoDLo}{0.10}
\newcommand{\AtwoDHi}{0.69}
\newcommand{\AtwoT}{2.37}
\newcommand{\AtwoDf}{102.5}
\newcommand{\AtwoP}{\pval{.010}}
\newcommand{\AtwoWP}{\pval{.069}}
\newcommand{\AtwoHL}{3.0}
\newcommand{\AtwoHLLo}{0.0}
\newcommand{\AtwoHLHi}{8.0}
\newcommand{\AtwoRB}{0.13}
\newcommand{\AtwoRBLo}{\ensuremath{-}0.05}
\newcommand{\AtwoRBHi}{0.30}
\newcommand{\AthreeN}{200}
\newcommand{\AthreeM}{18.1}
\newcommand{\AthreeLo}{16.9}
\newcommand{\AthreeHi}{19.2}
\newcommand{\AthreeDz}{2.19}
\newcommand{\AthreeDzLo}{1.93}
\newcommand{\AthreeDzHi}{2.44}
\newcommand{\AthreeT}{30.96}
\newcommand{\AthreeDf}{199}
\newcommand{\AthreeP}{\plt{.001}}
\newcommand{\AthreeWP}{\plt{.001}}
\newcommand{\AthreeHL}{17.5}
\newcommand{\AthreeHLLo}{17.5}
\newcommand{\AthreeHLHi}{18.5}
\newcommand{\AthreeSD}{8.3}
\newcommand{\AfourN}{190}
\newcommand{\AfourF}{2.33}
\newcommand{\AfourDfa}{5}
\newcommand{\AfourDfb}{58.8}
\newcommand{\AfourP}{\pval{.054}}
\newcommand{\AfourEta}{0.062}
\newcommand{\AfourEtaLo}{0.000}
\newcommand{\AfourEtaHi}{0.116}
\newcommand{\AfourOmega}{0.036}
\newcommand{\AfourH}{9.51}
\newcommand{\AfourKWdf}{5}
\newcommand{\AfourKWP}{\pval{.090}}
\newcommand{\AfourEps}{0.050}
\newcommand{\AfourBestPair}{Prague 1 minus Prague 7}
\newcommand{\AfourBestDiff}{8.5}
\newcommand{\AfourBestLo}{\ensuremath{-}0.2}
\newcommand{\AfourBestHi}{17.3}
\newcommand{\AfourBestP}{\pval{.062}}
\newcommand{\AfourNsigGH}{0}
\newcommand{\AfourNsigDunn}{0}
\newcommand{\AfourMax}{Prague 1}
\newcommand{\AfourMin}{Prague 3}
\newcommand{\AfourMaxM}{86.6}
\newcommand{\AfourMinM}{74.9}
\newcommand{\AfiveN}{199}
\newcommand{\AfiveChi}{5.31}
\newcommand{\AfiveDf}{2}
\newcommand{\AfiveP}{\pval{.070}}
\newcommand{\AfiveFisherP}{\pval{.074}}
\newcommand{\AfiveMinExp}{3.70}
\newcommand{\AfiveV}{0.16}
\newcommand{\AfiveVLo}{0.05}
\newcommand{\AfiveVHi}{0.31}
\newcommand{\AfiveShareInd}{33.6}
\newcommand{\AfiveShareLoc}{27.5}
\newcommand{\AfiveShareInt}{63.6}
\newcommand{\AfiveNInt}{11}
\newcommand{\AfiveKInt}{7}
\newcommand{\AfiveResInt}{2.16}
\newcommand{\AfiveResInd}{\ensuremath{-}0.04}
\newcommand{\AfiveResLoc}{\ensuremath{-}1.09}
\newcommand{\AsixN}{200}
\newcommand{\AsixR}{\ensuremath{-}0.17}
\newcommand{\AsixRLo}{\ensuremath{-}0.30}
\newcommand{\AsixRHi}{\ensuremath{-}0.03}
\newcommand{\AsixRsq}{2.7}
\newcommand{\AsixT}{\ensuremath{-}2.36}
\newcommand{\AsixDf}{198}
\newcommand{\AsixP}{\pval{.010}}
\newcommand{\AsixRho}{\ensuremath{-}0.18}
\newcommand{\AsixRhoLo}{\ensuremath{-}0.31}
\newcommand{\AsixRhoHi}{\ensuremath{-}0.03}
\newcommand{\AsixRhoP}{\pval{.006}}
\newcommand{\AsixTau}{\ensuremath{-}0.12}
\newcommand{\AsixTauLo}{\ensuremath{-}0.23}
\newcommand{\AsixTauHi}{\ensuremath{-}0.02}
\newcommand{\AsixTauP}{\pval{.006}}
\newcommand{\AsevenN}{200}
\newcommand{\AsevenB}{\ensuremath{-}1.85}
\newcommand{\AsevenLo}{\ensuremath{-}3.74}
\newcommand{\AsevenHi}{0.05}
\newcommand{\AsevenP}{\pval{.057}}
\newcommand{\AsevenT}{\ensuremath{-}1.92}
\newcommand{\AsevenDf}{198}
\newcommand{\AsevenRsq}{0.027}
\newcommand{\AsevenLogB}{\ensuremath{-}0.020}
\newcommand{\AsevenPct}{\ensuremath{-}2.0}
\newcommand{\AsevenPctLo}{\ensuremath{-}4.2}
\newcommand{\AsevenPctHi}{0.1}
\newcommand{\AsevenLogP}{\pval{.067}}
\newcommand{\AsevenLogT}{\ensuremath{-}1.84}
\newcommand{\AsevenLogRsq}{0.023}
\newcommand{\AsevenNinfl}{11}
\newcommand{\AsevenBinfl}{\ensuremath{-}1.21}
\newcommand{\AsevenMaxCook}{0.39}
\newcommand{\AsevenMaxDist}{8.1}
\newcommand{\AeightN}{199}
\newcommand{\AeightBa}{\ensuremath{-}1.84}
\newcommand{\AeightBb}{\ensuremath{-}1.77}
\newcommand{\AeightBbLo}{\ensuremath{-}3.58}
\newcommand{\AeightBbHi}{0.03}
\newcommand{\AeightBaLo}{\ensuremath{-}3.74}
\newcommand{\AeightBaHi}{0.06}
\newcommand{\AeightChange}{0.07}
\newcommand{\AeightRelChange}{3.5}
\newcommand{\AeightF}{6.67}
\newcommand{\AeightDfa}{2}
\newcommand{\AeightDfb}{195}
\newcommand{\AeightP}{\pval{.002}}
\newcommand{\AeightDeltaR}{0.062}
\newcommand{\AeightRsqA}{0.027}
\newcommand{\AeightRsqB}{0.090}
\newcommand{\AeightVifDist}{1.02}
\newcommand{\AeightVifType}{1.02}
\newcommand{\AeightGVIFadj}{1.01}
\newcommand{\AeightLocal}{6.1}
\newcommand{\AeightLocalLo}{1.9}
\newcommand{\AeightLocalHi}{10.3}
\newcommand{\AeightInt}{12.8}
\newcommand{\AeightIntLo}{\ensuremath{-}4.7}
\newcommand{\AeightIntHi}{30.4}
\newcommand{\AeightNinfl}{13}
\newcommand{\AeightBinfl}{\ensuremath{-}1.47}
\newcommand{\AnineN}{199}
\newcommand{\AnineNind}{137}
\newcommand{\AnineNch}{62}
\newcommand{\AnineInt}{\ensuremath{-}1.56}
\newcommand{\AnineIntLo}{\ensuremath{-}6.83}
\newcommand{\AnineIntHi}{3.70}
\newcommand{\AnineF}{1.00}
\newcommand{\AnineDfa}{1}
\newcommand{\AnineDfb}{195}
\newcommand{\AnineP}{\pval{.318}}
\newcommand{\AnineDeltaR}{0.005}
\newcommand{\AnineSlopeInd}{\ensuremath{-}1.33}
\newcommand{\AnineSlopeIndLo}{\ensuremath{-}2.95}
\newcommand{\AnineSlopeIndHi}{0.30}
\newcommand{\AnineSlopeCh}{\ensuremath{-}2.89}
\newcommand{\AnineSlopeChLo}{\ensuremath{-}7.90}
\newcommand{\AnineSlopeChHi}{2.12}
\newcommand{\AnineMaxDiff}{6.83}
\newcommand{\AtenN}{199}
\newcommand{\AtenEvents}{98}
\newcommand{\AtenAME}{\ensuremath{-}4.9}
\newcommand{\AtenAMELo}{\ensuremath{-}10.0}
\newcommand{\AtenAMEHi}{0.2}
\newcommand{\AtenOR}{0.81}
\newcommand{\AtenORLo}{0.65}
\newcommand{\AtenORHi}{1.02}
\newcommand{\AtenZ}{\ensuremath{-}1.81}
\newcommand{\AtenP}{\pval{.070}}
\newcommand{\AtenMcF}{0.039}
\newcommand{\AtenSampleMedian}{82.5}
\newcommand{\AtenShareHigh}{49.2}
\newcommand{\RobMinR}{\ensuremath{-}0.19}
\newcommand{\RobMaxR}{\ensuremath{-}0.12}
\newcommand{\ClnEmpty}{99}
\newcommand{\ClnDistrict}{20}
\newcommand{\ClnType}{1}
\newcommand{\ClnCap}{12}
\newcommand{\ClnEsp}{12}
\newcommand{\ClnSep}{6}
\newcommand{\ClnDup}{1}
\newcommand{\Nraw}{300}
\newcommand{\Nfilled}{201}
\newcommand{\Ntype}{199}
\newcommand{\PriceMean}{82.8}
\newcommand{\PriceSD}{14.8}
\newcommand{\PriceMedian}{82.5}
\newcommand{\PriceIQR}{15.0}
\newcommand{\PriceSkew}{0.53}
\newcommand{\PriceMin}{40}
\newcommand{\PriceMax}{140}
\newcommand{\DistSkew}{1.14}
\newcommand{\DistMedian}{1.64}
\newcommand{\DistMax}{8.1}
\newcommand{\DiffSkew}{0.64}
\newcommand{\AsevenSkewLin}{0.45}
\newcommand{\AsevenSkewLog}{\ensuremath{-}0.34}
\newcommand{\AoneDiffAbs}{7.2}
\newcommand{\AoneDir}{below}
\newcommand{\AonePctRel}{8.0}
\newcommand{\NstudentsCafes}{20}
\newcommand{\AthreePctEsp}{28}
\newcommand{\EspMean}{64.7}
\newcommand{\Ntests}{56}


\usetheme{default}
\usefonttheme{serif}
\setbeamertemplate{navigation symbols}{}
\setbeamercolor{background canvas}{bg=white}
\setbeamercolor{title}{fg=espresso}
\setbeamercolor{frametitle}{fg=espresso}
\setbeamercolor{normal text}{fg=espresso}
\setbeamercolor{itemize item}{fg=copper}
\setbeamerfont{frametitle}{series=\bfseries,size=\Large}
\setbeamertemplate{frametitle}{\vspace{0.4em}\insertframetitle}
\setbeamertemplate{footline}{\hfill\color{greyc}\small\insertframenumber\hspace{1em}\vspace{0.6em}}
\setbeamertemplate{itemize items}[circle]

\graphicspath{{figures/}}   % all graphs live in presentation/figures/  (A1_figure.png ... A10_figure.png)

% how the notes (the spoken script) are shown:
%   pdflatex slides.tex                                         -> slides only
%   pdflatex -jobname=slides_script "\def\withscript{1}%% Prague Cafe Project -- Batch 1 slides (Beamer). Numbers come from ../report/tables/numbers.tex,
%% which analysis/03_tests.R writes. Graphs are read from figures/ (A1_figure.png ... A10_figure.png),
%% which analysis/03_tests.R writes into this folder. Compile from this folder:  latexmk -pdf slides.tex
\documentclass[aspectratio=169,11pt]{beamer}
\usepackage[T1]{fontenc}
\usepackage{mathpazo}
\usepackage{graphicx}
\usepackage{booktabs}

\definecolor{espresso}{HTML}{2B1D16}
\definecolor{copper}{HTML}{B8622B}
\definecolor{pale}{HTML}{E7E2DC}
\definecolor{greyc}{HTML}{8A7F76}

% macros that numbers.tex expects
\newcommand{\pval}[1]{$p$\,=\,#1}
\newcommand{\plt}[1]{$p$\,<\,#1}
\newcommand{\CI}[2]{[#1,\,#2]}
% numbers.tex -- written by analysis/03_tests.R; do not edit by hand
\newcommand{\Nall}{200}
\newcommand{\AoneBench}{90}
\newcommand{\AoneN}{200}
\newcommand{\AoneMean}{82.8}
\newcommand{\AoneLo}{80.7}
\newcommand{\AoneHi}{84.8}
\newcommand{\AoneT}{\ensuremath{-}6.94}
\newcommand{\AoneDf}{199}
\newcommand{\AoneP}{\plt{.001}}
\newcommand{\AoneD}{\ensuremath{-}0.49}
\newcommand{\AoneDLo}{\ensuremath{-}0.64}
\newcommand{\AoneDHi}{\ensuremath{-}0.34}
\newcommand{\AoneWmed}{82.5}
\newcommand{\AoneWP}{\plt{.001}}
\newcommand{\AoneBLo}{80.8}
\newcommand{\AoneBHi}{84.8}
\newcommand{\AtwoNa}{67}
\newcommand{\AtwoNb}{133}
\newcommand{\AtwoMa}{86.6}
\newcommand{\AtwoMb}{80.8}
\newcommand{\AtwoDiff}{5.7}
\newcommand{\AtwoLo}{0.9}
\newcommand{\AtwoHi}{10.5}
\newcommand{\AtwoD}{0.39}
\newcommand{\AtwoDLo}{0.10}
\newcommand{\AtwoDHi}{0.69}
\newcommand{\AtwoT}{2.37}
\newcommand{\AtwoDf}{102.5}
\newcommand{\AtwoP}{\pval{.010}}
\newcommand{\AtwoWP}{\pval{.069}}
\newcommand{\AtwoHL}{3.0}
\newcommand{\AtwoHLLo}{0.0}
\newcommand{\AtwoHLHi}{8.0}
\newcommand{\AtwoRB}{0.13}
\newcommand{\AtwoRBLo}{\ensuremath{-}0.05}
\newcommand{\AtwoRBHi}{0.30}
\newcommand{\AthreeN}{200}
\newcommand{\AthreeM}{18.1}
\newcommand{\AthreeLo}{16.9}
\newcommand{\AthreeHi}{19.2}
\newcommand{\AthreeDz}{2.19}
\newcommand{\AthreeDzLo}{1.93}
\newcommand{\AthreeDzHi}{2.44}
\newcommand{\AthreeT}{30.96}
\newcommand{\AthreeDf}{199}
\newcommand{\AthreeP}{\plt{.001}}
\newcommand{\AthreeWP}{\plt{.001}}
\newcommand{\AthreeHL}{17.5}
\newcommand{\AthreeHLLo}{17.5}
\newcommand{\AthreeHLHi}{18.5}
\newcommand{\AthreeSD}{8.3}
\newcommand{\AfourN}{190}
\newcommand{\AfourF}{2.33}
\newcommand{\AfourDfa}{5}
\newcommand{\AfourDfb}{58.8}
\newcommand{\AfourP}{\pval{.054}}
\newcommand{\AfourEta}{0.062}
\newcommand{\AfourEtaLo}{0.000}
\newcommand{\AfourEtaHi}{0.116}
\newcommand{\AfourOmega}{0.036}
\newcommand{\AfourH}{9.51}
\newcommand{\AfourKWdf}{5}
\newcommand{\AfourKWP}{\pval{.090}}
\newcommand{\AfourEps}{0.050}
\newcommand{\AfourBestPair}{Prague 1 minus Prague 7}
\newcommand{\AfourBestDiff}{8.5}
\newcommand{\AfourBestLo}{\ensuremath{-}0.2}
\newcommand{\AfourBestHi}{17.3}
\newcommand{\AfourBestP}{\pval{.062}}
\newcommand{\AfourNsigGH}{0}
\newcommand{\AfourNsigDunn}{0}
\newcommand{\AfourMax}{Prague 1}
\newcommand{\AfourMin}{Prague 3}
\newcommand{\AfourMaxM}{86.6}
\newcommand{\AfourMinM}{74.9}
\newcommand{\AfiveN}{199}
\newcommand{\AfiveChi}{5.31}
\newcommand{\AfiveDf}{2}
\newcommand{\AfiveP}{\pval{.070}}
\newcommand{\AfiveFisherP}{\pval{.074}}
\newcommand{\AfiveMinExp}{3.70}
\newcommand{\AfiveV}{0.16}
\newcommand{\AfiveVLo}{0.05}
\newcommand{\AfiveVHi}{0.31}
\newcommand{\AfiveShareInd}{33.6}
\newcommand{\AfiveShareLoc}{27.5}
\newcommand{\AfiveShareInt}{63.6}
\newcommand{\AfiveNInt}{11}
\newcommand{\AfiveKInt}{7}
\newcommand{\AfiveResInt}{2.16}
\newcommand{\AfiveResInd}{\ensuremath{-}0.04}
\newcommand{\AfiveResLoc}{\ensuremath{-}1.09}
\newcommand{\AsixN}{200}
\newcommand{\AsixR}{\ensuremath{-}0.17}
\newcommand{\AsixRLo}{\ensuremath{-}0.30}
\newcommand{\AsixRHi}{\ensuremath{-}0.03}
\newcommand{\AsixRsq}{2.7}
\newcommand{\AsixT}{\ensuremath{-}2.36}
\newcommand{\AsixDf}{198}
\newcommand{\AsixP}{\pval{.010}}
\newcommand{\AsixRho}{\ensuremath{-}0.18}
\newcommand{\AsixRhoLo}{\ensuremath{-}0.31}
\newcommand{\AsixRhoHi}{\ensuremath{-}0.03}
\newcommand{\AsixRhoP}{\pval{.006}}
\newcommand{\AsixTau}{\ensuremath{-}0.12}
\newcommand{\AsixTauLo}{\ensuremath{-}0.23}
\newcommand{\AsixTauHi}{\ensuremath{-}0.02}
\newcommand{\AsixTauP}{\pval{.006}}
\newcommand{\AsevenN}{200}
\newcommand{\AsevenB}{\ensuremath{-}1.85}
\newcommand{\AsevenLo}{\ensuremath{-}3.74}
\newcommand{\AsevenHi}{0.05}
\newcommand{\AsevenP}{\pval{.057}}
\newcommand{\AsevenT}{\ensuremath{-}1.92}
\newcommand{\AsevenDf}{198}
\newcommand{\AsevenRsq}{0.027}
\newcommand{\AsevenLogB}{\ensuremath{-}0.020}
\newcommand{\AsevenPct}{\ensuremath{-}2.0}
\newcommand{\AsevenPctLo}{\ensuremath{-}4.2}
\newcommand{\AsevenPctHi}{0.1}
\newcommand{\AsevenLogP}{\pval{.067}}
\newcommand{\AsevenLogT}{\ensuremath{-}1.84}
\newcommand{\AsevenLogRsq}{0.023}
\newcommand{\AsevenNinfl}{11}
\newcommand{\AsevenBinfl}{\ensuremath{-}1.21}
\newcommand{\AsevenMaxCook}{0.39}
\newcommand{\AsevenMaxDist}{8.1}
\newcommand{\AeightN}{199}
\newcommand{\AeightBa}{\ensuremath{-}1.84}
\newcommand{\AeightBb}{\ensuremath{-}1.77}
\newcommand{\AeightBbLo}{\ensuremath{-}3.58}
\newcommand{\AeightBbHi}{0.03}
\newcommand{\AeightBaLo}{\ensuremath{-}3.74}
\newcommand{\AeightBaHi}{0.06}
\newcommand{\AeightChange}{0.07}
\newcommand{\AeightRelChange}{3.5}
\newcommand{\AeightF}{6.67}
\newcommand{\AeightDfa}{2}
\newcommand{\AeightDfb}{195}
\newcommand{\AeightP}{\pval{.002}}
\newcommand{\AeightDeltaR}{0.062}
\newcommand{\AeightRsqA}{0.027}
\newcommand{\AeightRsqB}{0.090}
\newcommand{\AeightVifDist}{1.02}
\newcommand{\AeightVifType}{1.02}
\newcommand{\AeightGVIFadj}{1.01}
\newcommand{\AeightLocal}{6.1}
\newcommand{\AeightLocalLo}{1.9}
\newcommand{\AeightLocalHi}{10.3}
\newcommand{\AeightInt}{12.8}
\newcommand{\AeightIntLo}{\ensuremath{-}4.7}
\newcommand{\AeightIntHi}{30.4}
\newcommand{\AeightNinfl}{13}
\newcommand{\AeightBinfl}{\ensuremath{-}1.47}
\newcommand{\AnineN}{199}
\newcommand{\AnineNind}{137}
\newcommand{\AnineNch}{62}
\newcommand{\AnineInt}{\ensuremath{-}1.56}
\newcommand{\AnineIntLo}{\ensuremath{-}6.83}
\newcommand{\AnineIntHi}{3.70}
\newcommand{\AnineF}{1.00}
\newcommand{\AnineDfa}{1}
\newcommand{\AnineDfb}{195}
\newcommand{\AnineP}{\pval{.318}}
\newcommand{\AnineDeltaR}{0.005}
\newcommand{\AnineSlopeInd}{\ensuremath{-}1.33}
\newcommand{\AnineSlopeIndLo}{\ensuremath{-}2.95}
\newcommand{\AnineSlopeIndHi}{0.30}
\newcommand{\AnineSlopeCh}{\ensuremath{-}2.89}
\newcommand{\AnineSlopeChLo}{\ensuremath{-}7.90}
\newcommand{\AnineSlopeChHi}{2.12}
\newcommand{\AnineMaxDiff}{6.83}
\newcommand{\AtenN}{199}
\newcommand{\AtenEvents}{98}
\newcommand{\AtenAME}{\ensuremath{-}4.9}
\newcommand{\AtenAMELo}{\ensuremath{-}10.0}
\newcommand{\AtenAMEHi}{0.2}
\newcommand{\AtenOR}{0.81}
\newcommand{\AtenORLo}{0.65}
\newcommand{\AtenORHi}{1.02}
\newcommand{\AtenZ}{\ensuremath{-}1.81}
\newcommand{\AtenP}{\pval{.070}}
\newcommand{\AtenMcF}{0.039}
\newcommand{\AtenSampleMedian}{82.5}
\newcommand{\AtenShareHigh}{49.2}
\newcommand{\RobMinR}{\ensuremath{-}0.19}
\newcommand{\RobMaxR}{\ensuremath{-}0.12}
\newcommand{\ClnEmpty}{99}
\newcommand{\ClnDistrict}{20}
\newcommand{\ClnType}{1}
\newcommand{\ClnCap}{12}
\newcommand{\ClnEsp}{12}
\newcommand{\ClnSep}{6}
\newcommand{\ClnDup}{1}
\newcommand{\Nraw}{300}
\newcommand{\Nfilled}{201}
\newcommand{\Ntype}{199}
\newcommand{\PriceMean}{82.8}
\newcommand{\PriceSD}{14.8}
\newcommand{\PriceMedian}{82.5}
\newcommand{\PriceIQR}{15.0}
\newcommand{\PriceSkew}{0.53}
\newcommand{\PriceMin}{40}
\newcommand{\PriceMax}{140}
\newcommand{\DistSkew}{1.14}
\newcommand{\DistMedian}{1.64}
\newcommand{\DistMax}{8.1}
\newcommand{\DiffSkew}{0.64}
\newcommand{\AsevenSkewLin}{0.45}
\newcommand{\AsevenSkewLog}{\ensuremath{-}0.34}
\newcommand{\AoneDiffAbs}{7.2}
\newcommand{\AoneDir}{below}
\newcommand{\AonePctRel}{8.0}
\newcommand{\NstudentsCafes}{20}
\newcommand{\AthreePctEsp}{28}
\newcommand{\EspMean}{64.7}
\newcommand{\Ntests}{56}


\usetheme{default}
\usefonttheme{serif}
\setbeamertemplate{navigation symbols}{}
\setbeamercolor{background canvas}{bg=white}
\setbeamercolor{title}{fg=espresso}
\setbeamercolor{frametitle}{fg=espresso}
\setbeamercolor{normal text}{fg=espresso}
\setbeamercolor{itemize item}{fg=copper}
\setbeamerfont{frametitle}{series=\bfseries,size=\Large}
\setbeamertemplate{frametitle}{\vspace{0.4em}\insertframetitle}
\setbeamertemplate{footline}{\hfill\color{greyc}\small\insertframenumber\hspace{1em}\vspace{0.6em}}
\setbeamertemplate{itemize items}[circle]

\graphicspath{{figures/}}   % all graphs live in presentation/figures/  (A1_figure.png ... A10_figure.png)

% how the notes (the spoken script) are shown:
%   pdflatex slides.tex                                         -> slides only
%   pdflatex -jobname=slides_script "\def\withscript{1}%% Prague Cafe Project -- Batch 1 slides (Beamer). Numbers come from ../report/tables/numbers.tex,
%% which analysis/03_tests.R writes. Graphs are read from figures/ (A1_figure.png ... A10_figure.png),
%% which analysis/03_tests.R writes into this folder. Compile from this folder:  latexmk -pdf slides.tex
\documentclass[aspectratio=169,11pt]{beamer}
\usepackage[T1]{fontenc}
\usepackage{mathpazo}
\usepackage{graphicx}
\usepackage{booktabs}

\definecolor{espresso}{HTML}{2B1D16}
\definecolor{copper}{HTML}{B8622B}
\definecolor{pale}{HTML}{E7E2DC}
\definecolor{greyc}{HTML}{8A7F76}

% macros that numbers.tex expects
\newcommand{\pval}[1]{$p$\,=\,#1}
\newcommand{\plt}[1]{$p$\,<\,#1}
\newcommand{\CI}[2]{[#1,\,#2]}
\input{../report/tables/numbers.tex}

\usetheme{default}
\usefonttheme{serif}
\setbeamertemplate{navigation symbols}{}
\setbeamercolor{background canvas}{bg=white}
\setbeamercolor{title}{fg=espresso}
\setbeamercolor{frametitle}{fg=espresso}
\setbeamercolor{normal text}{fg=espresso}
\setbeamercolor{itemize item}{fg=copper}
\setbeamerfont{frametitle}{series=\bfseries,size=\Large}
\setbeamertemplate{frametitle}{\vspace{0.4em}\insertframetitle}
\setbeamertemplate{footline}{\hfill\color{greyc}\small\insertframenumber\hspace{1em}\vspace{0.6em}}
\setbeamertemplate{itemize items}[circle]

\graphicspath{{figures/}}   % all graphs live in presentation/figures/  (A1_figure.png ... A10_figure.png)

% how the notes (the spoken script) are shown:
%   pdflatex slides.tex                                         -> slides only
%   pdflatex -jobname=slides_script "\def\withscript{1}\input{slides.tex}"   -> slides + script after each slide
\ifdefined\withscript
  \setbeameroption{show notes}
  \setbeamertemplate{note page}[plain]
  \setbeamerfont{note page}{size=\large}
\fi

% one slide per question (A1 ... A10), each built with \qframe; the spoken script is in the \note after it
% \qframe{title}{graph file}{stat 1 big}{stat 1 small}{stat 2 big}{stat 2 small}{test / reading line}{takeaway}
\newcommand{\stat}[2]{\colorbox{pale}{\parbox{0.9\linewidth}{\color{copper}\Large\bfseries #1\par\color{espresso}\small #2}}\par\vspace{0.5em}}
\newcommand{\qnote}[1]{\par\vspace{0.3em}{\small\color{greyc}#1}}
\newcommand{\tk}[1]{\vfill\noindent\colorbox{copper}{\parbox{\dimexpr\linewidth-2\fboxsep}{\color{white}\bfseries\ Takeaway: #1}}\vspace{0.2em}}
\newcommand{\qframe}[8]{%
\begin{frame}{#1}
\begin{columns}[T]
\column{0.60\textwidth}\centering\includegraphics[width=0.92\linewidth]{#2}
\column{0.37\textwidth}
\stat{#3}{#4}
\stat{#5}{#6}
\qnote{#7}
\end{columns}
\tk{#8}
\end{frame}}

\title{Does a Prague cappuccino get cheaper away from the centre?}
\author{Rishin Kumar}
\institute{Business Analytics, BANL 1100/1101}
\date{}

\begin{document}

% Title ----------------------------------------------------------
{\setbeamercolor{background canvas}{bg=espresso}\setbeamercolor{normal text}{fg=white}\setbeamercolor{itemize body}{fg=white}
\begin{frame}[plain]
\color{white}
\vspace{2.2cm}
{\Huge\bfseries Does a Prague cappuccino get cheaper away from the centre?\par}
\vspace{0.8cm}
{\large\color{copper!60!white} Price, distance and caf\'e type in \Nall{} caf\'es \,\textperiodcentered\, Batch 1\par}
\vspace{1.5cm}
{\small Rishin Kumar \,\textperiodcentered\, Business Analytics, BANL 1100/1101}
\end{frame}}
\note{Good morning, everyone. I'm Rishin. I looked at one question in the class's Prague caf\'e dataset: does a cappuccino get cheaper as you move away from the centre? I planned ten analyses in advance. For each finding, I'll tell you which test backs it up and why I trust it. The data: \Nall{} caf\'es after cleaning, visited by ten students. The average cappuccino costs \PriceMean{} crowns, and the typical caf\'e is \DistMedian{} kilometres from the Orloj. Everything I say is ``in this sample'': it is a convenience sample, not a random sample of Prague. Cleaning is all done in an R script, so anyone can rerun it and get my numbers.}

% One file per question ------------------------------------------
%% A1.tex -- slide for question A1; the graph is figures/A1_figure.png
\qframe{A1: Is the average price different from my prior guess?}{A1_figure.png}
  {\AoneMean{} CZK}{mean cappuccino, 95\% CI \CI{\AoneLo}{\AoneHi}}
  {$d$ = \AoneD}{one-sample $t$-test vs.\ \AoneBench{} CZK, \AoneP}
  {Wilcoxon and BCa bootstrap agree. Prior guess: \AoneBench{} CZK.}
  {A typical cappuccino costs about 83 CZK.}
\note{A1. What does a cappuccino cost? About \AoneMean{} crowns, with a 95 percent interval from \AoneLo{} to \AoneHi. [Say your own prior guess here, when you wrote it down, and the result. The file currently uses \AoneBench{} crowns as a placeholder.] I used a one-sample $t$-test. Defence: the price is only mildly skewed and the sample is large, and the Wilcoxon test and a bootstrap interval give the same answer.}

%% A2.tex -- slide for question A2; the graph is figures/A2_figure.png
\qframe{A2: Is Prague 1 dearer than the rest of Prague?}{A2_figure.png}
  {+\AtwoDiff{} CZK}{Prague 1 $-$ rest, 95\% CI \CI{\AtwoLo}{\AtwoHi}}
  {$d$ = \AtwoD}{Welch $t$-test, \AtwoP{} (one-tailed)}
  {Rank-sum check: \AtwoWP, shift \AtwoHL{} CZK. The tests only partly agree.}
  {Prague 1 is probably a few crowns dearer, but the tests only partly agree.}
\note{A2. Prague 1 averages \AtwoMa{} crowns against \AtwoMb{} elsewhere: \AtwoDiff{} crowns more, interval \AtwoLo{} to \AtwoHi, a small-to-medium effect. Welch's test, one-tailed as planned, gives \AtwoP. But the rank-based Wilcoxon test gives \AtwoWP{} and a typical shift of only \AtwoHL{} crowns. Defence: I chose Welch because the question is about averages and both groups are large, but I report the disagreement rather than hide it. A few very dear Prague 1 caf\'es pull the average up, so the honest summary is: probably 3 to 6 crowns dearer, under ten percent of the price.}

%% A3.tex -- slide for question A3; the graph is figures/A3_figure.png
\qframe{A3: How much more is a cappuccino than an espresso?}{A3_figure.png}
  {+\AthreeM{} CZK}{cappuccino $-$ espresso, 95\% CI \CI{\AthreeLo}{\AthreeHi}}
  {$d_z$ = \AthreeDz}{paired $t$-test, \AthreeP{} (one-tailed)}
  {About \AthreePctEsp\% on top of the espresso price. Wilcoxon agrees.}
  {A cappuccino costs about 28\% more than an espresso.}
\note{A3. In the same caf\'e, the cappuccino costs \AthreeM{} crowns more than the espresso, interval \AthreeLo{} to \AthreeHi: about \AthreePctEsp{} percent more. Paired $t$-test, one-tailed because I expected the cappuccino to be dearer, \AthreeP. Defence: pairing compares each caf\'e with itself, so differences between caf\'es cannot distort it, and the Wilcoxon test, which makes no normality assumption, agrees.}

%% A4.tex -- slide for question A4; the graph is figures/A4_figure.png
\qframe{A4: Do prices differ across the six districts?}{A4_figure.png}
  {$F(\AfourDfa,\,\AfourDfb)$ = \AfourF}{Welch ANOVA, \AfourP}
  {$\eta^2$ = \AfourEta}{95\% CI \CI{\AfourEtaLo}{\AfourEtaHi}}
  {Kruskal--Wallis agrees (\AfourKWP). No pair of districts differs after adjustment.}
  {No clear price differences between districts.}
\note{A4. Welch's ANOVA gives \AfourP{} and the Kruskal--Wallis test \AfourKWP. Districts explain only about 6 percent of the variation. No single pair of districts differs once I correct for comparing many pairs. Defence: with fifteen pairs some will look different by chance, which is why I used Games--Howell and Dunn--Holm corrections. The biggest gap, \AfourBestPair, is \AfourBestDiff{} crowns, but its interval includes zero.}

%% A5.tex -- slide for question A5; the graph is figures/A5_figure.png
\qframe{A5: Is caf\'e type linked to being in Prague 1?}{A5_figure.png}
  {\AfiveFisherP}{Fisher's exact test}
  {$V$ = \AfiveV}{Cram\'er's $V$, 95\% CI \CI{\AfiveVLo}{\AfiveVHi}}
  {\AfiveKInt{} of \AfiveNInt{} international chains are in Prague 1.}
  {International chains cluster in Prague 1; suggestive, not conclusive.}
\note{A5. \AfiveKInt{} of the \AfiveNInt{} international chains, about \AfiveShareInt{} percent, are in Prague 1, compared with about \AfiveShareInd{} percent of independents. Fisher's exact test gives \AfiveFisherP. Defence: I used Fisher's test instead of chi-square because one expected count is small. The pattern is suggestive, not conclusive.}

%% A6.tex -- slide for question A6; the graph is figures/A6_figure.png
\qframe{A6: Does price fall with distance from the Orloj?}{A6_figure.png}
  {$r$ = \AsixR}{Pearson, 95\% CI \CI{\AsixRLo}{\AsixRHi}}
  {\AsixP}{one-tailed; distance shares \AsixRsq\% of price variance}
  {Spearman $\rho$ = \AsixRho{} and Kendall $\tau$ = \AsixTau{} agree.}
  {Farther from the centre, prices are slightly lower.}
\note{A6. The scatterplot shows a downward trend, but a weak one. Pearson's correlation is \AsixR, one-tailed \AsixP. Distance explains under three percent of the variation in price. Defence: I decided the direction in advance, which is why a one-tailed test is fair. Distance is skewed, so I also ran Spearman and Kendall correlations, \AsixRho{} and \AsixTau: the same story.}

%% A7.tex -- slide for question A7; the graph is figures/A7_figure.png
\qframe{A7: By how much does price change per kilometre?}{A7_figure.png}
  {\AsevenPct\% per km}{log-linear, robust 95\% CI \CI{\AsevenPctLo}{\AsevenPctHi}}
  {\AsevenLogP}{two-tailed, HC3 robust standard errors}
  {About \AsevenB{} CZK per km. \AsevenNinfl{} influential caf\'es: without them \AsevenBinfl{} CZK per km.}
  {About 2\% cheaper per km: a small, uncertain effect.}
\note{A7. Each extra kilometre goes with a price change of about \AsevenPct{} percent, roughly \AsevenB{} crowns. The interval runs from \AsevenPctLo{} to \AsevenPctHi{} percent, so it just touches zero. Defence: this test is two-tailed with robust standard errors, so it is stricter than A6, which is why it is less conclusive. Eleven influential caf\'es matter: without them the slope is about \AsevenBinfl{} crowns, still negative. This shows an association only: rent, tourist traffic and caf\'e type all move with distance, so the data cannot say why prices fall.}

%% A8.tex -- slide for question A8; the graph is figures/A8_figure.png
\qframe{A8: Does the distance effect survive controlling for caf\'e type?}{A8_figure.png}
  {\AeightBa{} $\to$ \AeightBb{}}{distance slope (CZK/km) without $\to$ with caf\'e type}
  {$\Delta R^2$ = \AeightDeltaR}{nested $F$-test, \AeightP}
  {Local chains charge +\AeightLocal{} CZK over independents at the same distance.}
  {Distance still matters after caf\'e type, and chains charge more.}
\note{A8. Once I add caf\'e type, the distance slope barely moves, from \AeightBa{} to \AeightBb{} crowns per kilometre. Caf\'e type itself improves the model, nested $F$-test \AeightP, and local chains charge about \AeightLocal{} crowns more than independents at the same distance. Defence: the variance inflation factors are close to one, so distance and type carry separate information.}

%% A9.tex -- slide for question A9; the graph is figures/A9_figure.png
\qframe{A9: Does the distance slope differ for independents and chains?}{A9_figure.png}
  {\AnineInt{} CZK/km}{interaction, robust 95\% CI \CI{\AnineIntLo}{\AnineIntHi}}
  {\AnineP}{nested $F$-test; $\Delta R^2$ = \AnineDeltaR}
  {Slopes: independents \AnineSlopeInd, chains \AnineSlopeCh{} CZK/km. No clear difference.}
  {No clear evidence that the distance effect differs for chains.}
\note{A9. The slope is \AnineSlopeInd{} crowns per kilometre for independents and \AnineSlopeCh{} for chains. That looks different, but the interaction has \AnineP{} and an interval from \AnineIntLo{} to \AnineIntHi. Defence: no clear difference is not the same as evidence the slopes are equal. The data are simply too noisy to tell.}

%% A10.tex -- slide for question A10; the graph is figures/A10_figure.png
\qframe{A10: How does the chance of an above-median price change per km?}{A10_figure.png}
  {\AtenAME{} pp per km}{average marginal effect, 95\% CI \CI{\AtenAMELo}{\AtenAMEHi}}
  {OR = \AtenOR}{logit, \AtenP}
  {Outcome: price above 83 CZK (the brief's median).}
  {Each extra km lowers the chance of an above-median price by about 5 points.}
\note{A10. I turned price into a yes-or-no question: is a cappuccino above 83 crowns, the median in the project brief? In a logistic regression each extra kilometre changes that probability by an average of \AtenAME{} percentage points, interval \AtenAMELo{} to \AtenAMEHi, \AtenP; the odds ratio is \AtenOR. Defence: I report percentage points because they are easier to read than odds. It points the same way as everything else, but the interval includes zero, so I am not overselling it. To close: I ran \Ntests{} tests in total, so one borderline $p$-value is not a finding, which is why I focus on the ten planned analyses and on effect sizes and intervals. This is a convenience sample, the data are observational, and caf\'es in the same district are not fully independent. Thank you; I am happy to take questions.}


\end{document}
"   -> slides + script after each slide
\ifdefined\withscript
  \setbeameroption{show notes}
  \setbeamertemplate{note page}[plain]
  \setbeamerfont{note page}{size=\large}
\fi

% one slide per question (A1 ... A10), each built with \qframe; the spoken script is in the \note after it
% \qframe{title}{graph file}{stat 1 big}{stat 1 small}{stat 2 big}{stat 2 small}{test / reading line}{takeaway}
\newcommand{\stat}[2]{\colorbox{pale}{\parbox{0.9\linewidth}{\color{copper}\Large\bfseries #1\par\color{espresso}\small #2}}\par\vspace{0.5em}}
\newcommand{\qnote}[1]{\par\vspace{0.3em}{\small\color{greyc}#1}}
\newcommand{\tk}[1]{\vfill\noindent\colorbox{copper}{\parbox{\dimexpr\linewidth-2\fboxsep}{\color{white}\bfseries\ Takeaway: #1}}\vspace{0.2em}}
\newcommand{\qframe}[8]{%
\begin{frame}{#1}
\begin{columns}[T]
\column{0.60\textwidth}\centering\includegraphics[width=0.92\linewidth]{#2}
\column{0.37\textwidth}
\stat{#3}{#4}
\stat{#5}{#6}
\qnote{#7}
\end{columns}
\tk{#8}
\end{frame}}

\title{Does a Prague cappuccino get cheaper away from the centre?}
\author{Rishin Kumar}
\institute{Business Analytics, BANL 1100/1101}
\date{}

\begin{document}

% Title ----------------------------------------------------------
{\setbeamercolor{background canvas}{bg=espresso}\setbeamercolor{normal text}{fg=white}\setbeamercolor{itemize body}{fg=white}
\begin{frame}[plain]
\color{white}
\vspace{2.2cm}
{\Huge\bfseries Does a Prague cappuccino get cheaper away from the centre?\par}
\vspace{0.8cm}
{\large\color{copper!60!white} Price, distance and caf\'e type in \Nall{} caf\'es \,\textperiodcentered\, Batch 1\par}
\vspace{1.5cm}
{\small Rishin Kumar \,\textperiodcentered\, Business Analytics, BANL 1100/1101}
\end{frame}}
\note{Good morning, everyone. I'm Rishin. I looked at one question in the class's Prague caf\'e dataset: does a cappuccino get cheaper as you move away from the centre? I planned ten analyses in advance. For each finding, I'll tell you which test backs it up and why I trust it. The data: \Nall{} caf\'es after cleaning, visited by ten students. The average cappuccino costs \PriceMean{} crowns, and the typical caf\'e is \DistMedian{} kilometres from the Orloj. Everything I say is ``in this sample'': it is a convenience sample, not a random sample of Prague. Cleaning is all done in an R script, so anyone can rerun it and get my numbers.}

% One file per question ------------------------------------------
%% A1.tex -- slide for question A1; the graph is figures/A1_figure.png
\qframe{A1: Is the average price different from my prior guess?}{A1_figure.png}
  {\AoneMean{} CZK}{mean cappuccino, 95\% CI \CI{\AoneLo}{\AoneHi}}
  {$d$ = \AoneD}{one-sample $t$-test vs.\ \AoneBench{} CZK, \AoneP}
  {Wilcoxon and BCa bootstrap agree. Prior guess: \AoneBench{} CZK.}
  {A typical cappuccino costs about 83 CZK.}
\note{A1. What does a cappuccino cost? About \AoneMean{} crowns, with a 95 percent interval from \AoneLo{} to \AoneHi. [Say your own prior guess here, when you wrote it down, and the result. The file currently uses \AoneBench{} crowns as a placeholder.] I used a one-sample $t$-test. Defence: the price is only mildly skewed and the sample is large, and the Wilcoxon test and a bootstrap interval give the same answer.}

%% A2.tex -- slide for question A2; the graph is figures/A2_figure.png
\qframe{A2: Is Prague 1 dearer than the rest of Prague?}{A2_figure.png}
  {+\AtwoDiff{} CZK}{Prague 1 $-$ rest, 95\% CI \CI{\AtwoLo}{\AtwoHi}}
  {$d$ = \AtwoD}{Welch $t$-test, \AtwoP{} (one-tailed)}
  {Rank-sum check: \AtwoWP, shift \AtwoHL{} CZK. The tests only partly agree.}
  {Prague 1 is probably a few crowns dearer, but the tests only partly agree.}
\note{A2. Prague 1 averages \AtwoMa{} crowns against \AtwoMb{} elsewhere: \AtwoDiff{} crowns more, interval \AtwoLo{} to \AtwoHi, a small-to-medium effect. Welch's test, one-tailed as planned, gives \AtwoP. But the rank-based Wilcoxon test gives \AtwoWP{} and a typical shift of only \AtwoHL{} crowns. Defence: I chose Welch because the question is about averages and both groups are large, but I report the disagreement rather than hide it. A few very dear Prague 1 caf\'es pull the average up, so the honest summary is: probably 3 to 6 crowns dearer, under ten percent of the price.}

%% A3.tex -- slide for question A3; the graph is figures/A3_figure.png
\qframe{A3: How much more is a cappuccino than an espresso?}{A3_figure.png}
  {+\AthreeM{} CZK}{cappuccino $-$ espresso, 95\% CI \CI{\AthreeLo}{\AthreeHi}}
  {$d_z$ = \AthreeDz}{paired $t$-test, \AthreeP{} (one-tailed)}
  {About \AthreePctEsp\% on top of the espresso price. Wilcoxon agrees.}
  {A cappuccino costs about 28\% more than an espresso.}
\note{A3. In the same caf\'e, the cappuccino costs \AthreeM{} crowns more than the espresso, interval \AthreeLo{} to \AthreeHi: about \AthreePctEsp{} percent more. Paired $t$-test, one-tailed because I expected the cappuccino to be dearer, \AthreeP. Defence: pairing compares each caf\'e with itself, so differences between caf\'es cannot distort it, and the Wilcoxon test, which makes no normality assumption, agrees.}

%% A4.tex -- slide for question A4; the graph is figures/A4_figure.png
\qframe{A4: Do prices differ across the six districts?}{A4_figure.png}
  {$F(\AfourDfa,\,\AfourDfb)$ = \AfourF}{Welch ANOVA, \AfourP}
  {$\eta^2$ = \AfourEta}{95\% CI \CI{\AfourEtaLo}{\AfourEtaHi}}
  {Kruskal--Wallis agrees (\AfourKWP). No pair of districts differs after adjustment.}
  {No clear price differences between districts.}
\note{A4. Welch's ANOVA gives \AfourP{} and the Kruskal--Wallis test \AfourKWP. Districts explain only about 6 percent of the variation. No single pair of districts differs once I correct for comparing many pairs. Defence: with fifteen pairs some will look different by chance, which is why I used Games--Howell and Dunn--Holm corrections. The biggest gap, \AfourBestPair, is \AfourBestDiff{} crowns, but its interval includes zero.}

%% A5.tex -- slide for question A5; the graph is figures/A5_figure.png
\qframe{A5: Is caf\'e type linked to being in Prague 1?}{A5_figure.png}
  {\AfiveFisherP}{Fisher's exact test}
  {$V$ = \AfiveV}{Cram\'er's $V$, 95\% CI \CI{\AfiveVLo}{\AfiveVHi}}
  {\AfiveKInt{} of \AfiveNInt{} international chains are in Prague 1.}
  {International chains cluster in Prague 1; suggestive, not conclusive.}
\note{A5. \AfiveKInt{} of the \AfiveNInt{} international chains, about \AfiveShareInt{} percent, are in Prague 1, compared with about \AfiveShareInd{} percent of independents. Fisher's exact test gives \AfiveFisherP. Defence: I used Fisher's test instead of chi-square because one expected count is small. The pattern is suggestive, not conclusive.}

%% A6.tex -- slide for question A6; the graph is figures/A6_figure.png
\qframe{A6: Does price fall with distance from the Orloj?}{A6_figure.png}
  {$r$ = \AsixR}{Pearson, 95\% CI \CI{\AsixRLo}{\AsixRHi}}
  {\AsixP}{one-tailed; distance shares \AsixRsq\% of price variance}
  {Spearman $\rho$ = \AsixRho{} and Kendall $\tau$ = \AsixTau{} agree.}
  {Farther from the centre, prices are slightly lower.}
\note{A6. The scatterplot shows a downward trend, but a weak one. Pearson's correlation is \AsixR, one-tailed \AsixP. Distance explains under three percent of the variation in price. Defence: I decided the direction in advance, which is why a one-tailed test is fair. Distance is skewed, so I also ran Spearman and Kendall correlations, \AsixRho{} and \AsixTau: the same story.}

%% A7.tex -- slide for question A7; the graph is figures/A7_figure.png
\qframe{A7: By how much does price change per kilometre?}{A7_figure.png}
  {\AsevenPct\% per km}{log-linear, robust 95\% CI \CI{\AsevenPctLo}{\AsevenPctHi}}
  {\AsevenLogP}{two-tailed, HC3 robust standard errors}
  {About \AsevenB{} CZK per km. \AsevenNinfl{} influential caf\'es: without them \AsevenBinfl{} CZK per km.}
  {About 2\% cheaper per km: a small, uncertain effect.}
\note{A7. Each extra kilometre goes with a price change of about \AsevenPct{} percent, roughly \AsevenB{} crowns. The interval runs from \AsevenPctLo{} to \AsevenPctHi{} percent, so it just touches zero. Defence: this test is two-tailed with robust standard errors, so it is stricter than A6, which is why it is less conclusive. Eleven influential caf\'es matter: without them the slope is about \AsevenBinfl{} crowns, still negative. This shows an association only: rent, tourist traffic and caf\'e type all move with distance, so the data cannot say why prices fall.}

%% A8.tex -- slide for question A8; the graph is figures/A8_figure.png
\qframe{A8: Does the distance effect survive controlling for caf\'e type?}{A8_figure.png}
  {\AeightBa{} $\to$ \AeightBb{}}{distance slope (CZK/km) without $\to$ with caf\'e type}
  {$\Delta R^2$ = \AeightDeltaR}{nested $F$-test, \AeightP}
  {Local chains charge +\AeightLocal{} CZK over independents at the same distance.}
  {Distance still matters after caf\'e type, and chains charge more.}
\note{A8. Once I add caf\'e type, the distance slope barely moves, from \AeightBa{} to \AeightBb{} crowns per kilometre. Caf\'e type itself improves the model, nested $F$-test \AeightP, and local chains charge about \AeightLocal{} crowns more than independents at the same distance. Defence: the variance inflation factors are close to one, so distance and type carry separate information.}

%% A9.tex -- slide for question A9; the graph is figures/A9_figure.png
\qframe{A9: Does the distance slope differ for independents and chains?}{A9_figure.png}
  {\AnineInt{} CZK/km}{interaction, robust 95\% CI \CI{\AnineIntLo}{\AnineIntHi}}
  {\AnineP}{nested $F$-test; $\Delta R^2$ = \AnineDeltaR}
  {Slopes: independents \AnineSlopeInd, chains \AnineSlopeCh{} CZK/km. No clear difference.}
  {No clear evidence that the distance effect differs for chains.}
\note{A9. The slope is \AnineSlopeInd{} crowns per kilometre for independents and \AnineSlopeCh{} for chains. That looks different, but the interaction has \AnineP{} and an interval from \AnineIntLo{} to \AnineIntHi. Defence: no clear difference is not the same as evidence the slopes are equal. The data are simply too noisy to tell.}

%% A10.tex -- slide for question A10; the graph is figures/A10_figure.png
\qframe{A10: How does the chance of an above-median price change per km?}{A10_figure.png}
  {\AtenAME{} pp per km}{average marginal effect, 95\% CI \CI{\AtenAMELo}{\AtenAMEHi}}
  {OR = \AtenOR}{logit, \AtenP}
  {Outcome: price above 83 CZK (the brief's median).}
  {Each extra km lowers the chance of an above-median price by about 5 points.}
\note{A10. I turned price into a yes-or-no question: is a cappuccino above 83 crowns, the median in the project brief? In a logistic regression each extra kilometre changes that probability by an average of \AtenAME{} percentage points, interval \AtenAMELo{} to \AtenAMEHi, \AtenP; the odds ratio is \AtenOR. Defence: I report percentage points because they are easier to read than odds. It points the same way as everything else, but the interval includes zero, so I am not overselling it. To close: I ran \Ntests{} tests in total, so one borderline $p$-value is not a finding, which is why I focus on the ten planned analyses and on effect sizes and intervals. This is a convenience sample, the data are observational, and caf\'es in the same district are not fully independent. Thank you; I am happy to take questions.}


\end{document}
"   -> slides + script after each slide
\ifdefined\withscript
  \setbeameroption{show notes}
  \setbeamertemplate{note page}[plain]
  \setbeamerfont{note page}{size=\large}
\fi

% one slide per question (A1 ... A10), each built with \qframe; the spoken script is in the \note after it
% \qframe{title}{graph file}{stat 1 big}{stat 1 small}{stat 2 big}{stat 2 small}{test / reading line}{takeaway}
\newcommand{\stat}[2]{\colorbox{pale}{\parbox{0.9\linewidth}{\color{copper}\Large\bfseries #1\par\color{espresso}\small #2}}\par\vspace{0.5em}}
\newcommand{\qnote}[1]{\par\vspace{0.3em}{\small\color{greyc}#1}}
\newcommand{\tk}[1]{\vfill\noindent\colorbox{copper}{\parbox{\dimexpr\linewidth-2\fboxsep}{\color{white}\bfseries\ Takeaway: #1}}\vspace{0.2em}}
\newcommand{\qframe}[8]{%
\begin{frame}{#1}
\begin{columns}[T]
\column{0.60\textwidth}\centering\includegraphics[width=0.92\linewidth]{#2}
\column{0.37\textwidth}
\stat{#3}{#4}
\stat{#5}{#6}
\qnote{#7}
\end{columns}
\tk{#8}
\end{frame}}

\title{Does a Prague cappuccino get cheaper away from the centre?}
\author{Rishin Kumar}
\institute{Business Analytics, BANL 1100/1101}
\date{}

\begin{document}

% Title ----------------------------------------------------------
{\setbeamercolor{background canvas}{bg=espresso}\setbeamercolor{normal text}{fg=white}\setbeamercolor{itemize body}{fg=white}
\begin{frame}[plain]
\color{white}
\vspace{2.2cm}
{\Huge\bfseries Does a Prague cappuccino get cheaper away from the centre?\par}
\vspace{0.8cm}
{\large\color{copper!60!white} Price, distance and caf\'e type in \Nall{} caf\'es \,\textperiodcentered\, Batch 1\par}
\vspace{1.5cm}
{\small Rishin Kumar \,\textperiodcentered\, Business Analytics, BANL 1100/1101}
\end{frame}}
\note{Good morning, everyone. I'm Rishin. I looked at one question in the class's Prague caf\'e dataset: does a cappuccino get cheaper as you move away from the centre? I planned ten analyses in advance. For each finding, I'll tell you which test backs it up and why I trust it. The data: \Nall{} caf\'es after cleaning, visited by ten students. The average cappuccino costs \PriceMean{} crowns, and the typical caf\'e is \DistMedian{} kilometres from the Orloj. Everything I say is ``in this sample'': it is a convenience sample, not a random sample of Prague. Cleaning is all done in an R script, so anyone can rerun it and get my numbers.}

% One file per question ------------------------------------------
%% A1.tex -- slide for question A1; the graph is figures/A1_figure.png
\qframe{A1: Is the average price different from my prior guess?}{A1_figure.png}
  {\AoneMean{} CZK}{mean cappuccino, 95\% CI \CI{\AoneLo}{\AoneHi}}
  {$d$ = \AoneD}{one-sample $t$-test vs.\ \AoneBench{} CZK, \AoneP}
  {Wilcoxon and BCa bootstrap agree. Prior guess: \AoneBench{} CZK.}
  {A typical cappuccino costs about 83 CZK.}
\note{A1. What does a cappuccino cost? About \AoneMean{} crowns, with a 95 percent interval from \AoneLo{} to \AoneHi. [Say your own prior guess here, when you wrote it down, and the result. The file currently uses \AoneBench{} crowns as a placeholder.] I used a one-sample $t$-test. Defence: the price is only mildly skewed and the sample is large, and the Wilcoxon test and a bootstrap interval give the same answer.}

%% A2.tex -- slide for question A2; the graph is figures/A2_figure.png
\qframe{A2: Is Prague 1 dearer than the rest of Prague?}{A2_figure.png}
  {+\AtwoDiff{} CZK}{Prague 1 $-$ rest, 95\% CI \CI{\AtwoLo}{\AtwoHi}}
  {$d$ = \AtwoD}{Welch $t$-test, \AtwoP{} (one-tailed)}
  {Rank-sum check: \AtwoWP, shift \AtwoHL{} CZK. The tests only partly agree.}
  {Prague 1 is probably a few crowns dearer, but the tests only partly agree.}
\note{A2. Prague 1 averages \AtwoMa{} crowns against \AtwoMb{} elsewhere: \AtwoDiff{} crowns more, interval \AtwoLo{} to \AtwoHi, a small-to-medium effect. Welch's test, one-tailed as planned, gives \AtwoP. But the rank-based Wilcoxon test gives \AtwoWP{} and a typical shift of only \AtwoHL{} crowns. Defence: I chose Welch because the question is about averages and both groups are large, but I report the disagreement rather than hide it. A few very dear Prague 1 caf\'es pull the average up, so the honest summary is: probably 3 to 6 crowns dearer, under ten percent of the price.}

%% A3.tex -- slide for question A3; the graph is figures/A3_figure.png
\qframe{A3: How much more is a cappuccino than an espresso?}{A3_figure.png}
  {+\AthreeM{} CZK}{cappuccino $-$ espresso, 95\% CI \CI{\AthreeLo}{\AthreeHi}}
  {$d_z$ = \AthreeDz}{paired $t$-test, \AthreeP{} (one-tailed)}
  {About \AthreePctEsp\% on top of the espresso price. Wilcoxon agrees.}
  {A cappuccino costs about 28\% more than an espresso.}
\note{A3. In the same caf\'e, the cappuccino costs \AthreeM{} crowns more than the espresso, interval \AthreeLo{} to \AthreeHi: about \AthreePctEsp{} percent more. Paired $t$-test, one-tailed because I expected the cappuccino to be dearer, \AthreeP. Defence: pairing compares each caf\'e with itself, so differences between caf\'es cannot distort it, and the Wilcoxon test, which makes no normality assumption, agrees.}

%% A4.tex -- slide for question A4; the graph is figures/A4_figure.png
\qframe{A4: Do prices differ across the six districts?}{A4_figure.png}
  {$F(\AfourDfa,\,\AfourDfb)$ = \AfourF}{Welch ANOVA, \AfourP}
  {$\eta^2$ = \AfourEta}{95\% CI \CI{\AfourEtaLo}{\AfourEtaHi}}
  {Kruskal--Wallis agrees (\AfourKWP). No pair of districts differs after adjustment.}
  {No clear price differences between districts.}
\note{A4. Welch's ANOVA gives \AfourP{} and the Kruskal--Wallis test \AfourKWP. Districts explain only about 6 percent of the variation. No single pair of districts differs once I correct for comparing many pairs. Defence: with fifteen pairs some will look different by chance, which is why I used Games--Howell and Dunn--Holm corrections. The biggest gap, \AfourBestPair, is \AfourBestDiff{} crowns, but its interval includes zero.}

%% A5.tex -- slide for question A5; the graph is figures/A5_figure.png
\qframe{A5: Is caf\'e type linked to being in Prague 1?}{A5_figure.png}
  {\AfiveFisherP}{Fisher's exact test}
  {$V$ = \AfiveV}{Cram\'er's $V$, 95\% CI \CI{\AfiveVLo}{\AfiveVHi}}
  {\AfiveKInt{} of \AfiveNInt{} international chains are in Prague 1.}
  {International chains cluster in Prague 1; suggestive, not conclusive.}
\note{A5. \AfiveKInt{} of the \AfiveNInt{} international chains, about \AfiveShareInt{} percent, are in Prague 1, compared with about \AfiveShareInd{} percent of independents. Fisher's exact test gives \AfiveFisherP. Defence: I used Fisher's test instead of chi-square because one expected count is small. The pattern is suggestive, not conclusive.}

%% A6.tex -- slide for question A6; the graph is figures/A6_figure.png
\qframe{A6: Does price fall with distance from the Orloj?}{A6_figure.png}
  {$r$ = \AsixR}{Pearson, 95\% CI \CI{\AsixRLo}{\AsixRHi}}
  {\AsixP}{one-tailed; distance shares \AsixRsq\% of price variance}
  {Spearman $\rho$ = \AsixRho{} and Kendall $\tau$ = \AsixTau{} agree.}
  {Farther from the centre, prices are slightly lower.}
\note{A6. The scatterplot shows a downward trend, but a weak one. Pearson's correlation is \AsixR, one-tailed \AsixP. Distance explains under three percent of the variation in price. Defence: I decided the direction in advance, which is why a one-tailed test is fair. Distance is skewed, so I also ran Spearman and Kendall correlations, \AsixRho{} and \AsixTau: the same story.}

%% A7.tex -- slide for question A7; the graph is figures/A7_figure.png
\qframe{A7: By how much does price change per kilometre?}{A7_figure.png}
  {\AsevenPct\% per km}{log-linear, robust 95\% CI \CI{\AsevenPctLo}{\AsevenPctHi}}
  {\AsevenLogP}{two-tailed, HC3 robust standard errors}
  {About \AsevenB{} CZK per km. \AsevenNinfl{} influential caf\'es: without them \AsevenBinfl{} CZK per km.}
  {About 2\% cheaper per km: a small, uncertain effect.}
\note{A7. Each extra kilometre goes with a price change of about \AsevenPct{} percent, roughly \AsevenB{} crowns. The interval runs from \AsevenPctLo{} to \AsevenPctHi{} percent, so it just touches zero. Defence: this test is two-tailed with robust standard errors, so it is stricter than A6, which is why it is less conclusive. Eleven influential caf\'es matter: without them the slope is about \AsevenBinfl{} crowns, still negative. This shows an association only: rent, tourist traffic and caf\'e type all move with distance, so the data cannot say why prices fall.}

%% A8.tex -- slide for question A8; the graph is figures/A8_figure.png
\qframe{A8: Does the distance effect survive controlling for caf\'e type?}{A8_figure.png}
  {\AeightBa{} $\to$ \AeightBb{}}{distance slope (CZK/km) without $\to$ with caf\'e type}
  {$\Delta R^2$ = \AeightDeltaR}{nested $F$-test, \AeightP}
  {Local chains charge +\AeightLocal{} CZK over independents at the same distance.}
  {Distance still matters after caf\'e type, and chains charge more.}
\note{A8. Once I add caf\'e type, the distance slope barely moves, from \AeightBa{} to \AeightBb{} crowns per kilometre. Caf\'e type itself improves the model, nested $F$-test \AeightP, and local chains charge about \AeightLocal{} crowns more than independents at the same distance. Defence: the variance inflation factors are close to one, so distance and type carry separate information.}

%% A9.tex -- slide for question A9; the graph is figures/A9_figure.png
\qframe{A9: Does the distance slope differ for independents and chains?}{A9_figure.png}
  {\AnineInt{} CZK/km}{interaction, robust 95\% CI \CI{\AnineIntLo}{\AnineIntHi}}
  {\AnineP}{nested $F$-test; $\Delta R^2$ = \AnineDeltaR}
  {Slopes: independents \AnineSlopeInd, chains \AnineSlopeCh{} CZK/km. No clear difference.}
  {No clear evidence that the distance effect differs for chains.}
\note{A9. The slope is \AnineSlopeInd{} crowns per kilometre for independents and \AnineSlopeCh{} for chains. That looks different, but the interaction has \AnineP{} and an interval from \AnineIntLo{} to \AnineIntHi. Defence: no clear difference is not the same as evidence the slopes are equal. The data are simply too noisy to tell.}

%% A10.tex -- slide for question A10; the graph is figures/A10_figure.png
\qframe{A10: How does the chance of an above-median price change per km?}{A10_figure.png}
  {\AtenAME{} pp per km}{average marginal effect, 95\% CI \CI{\AtenAMELo}{\AtenAMEHi}}
  {OR = \AtenOR}{logit, \AtenP}
  {Outcome: price above 83 CZK (the brief's median).}
  {Each extra km lowers the chance of an above-median price by about 5 points.}
\note{A10. I turned price into a yes-or-no question: is a cappuccino above 83 crowns, the median in the project brief? In a logistic regression each extra kilometre changes that probability by an average of \AtenAME{} percentage points, interval \AtenAMELo{} to \AtenAMEHi, \AtenP; the odds ratio is \AtenOR. Defence: I report percentage points because they are easier to read than odds. It points the same way as everything else, but the interval includes zero, so I am not overselling it. To close: I ran \Ntests{} tests in total, so one borderline $p$-value is not a finding, which is why I focus on the ten planned analyses and on effect sizes and intervals. This is a convenience sample, the data are observational, and caf\'es in the same district are not fully independent. Thank you; I am happy to take questions.}


\end{document}
"   -> slides + script after each slide
\ifdefined\withscript
  \setbeameroption{show notes}
  \setbeamertemplate{note page}[plain]
  \setbeamerfont{note page}{size=\large}
\fi

% one slide per question (A1 ... A10), each built with \qframe; the spoken script is in the \note after it
% \qframe{title}{graph file}{stat 1 big}{stat 1 small}{stat 2 big}{stat 2 small}{test / reading line}{takeaway}
\newcommand{\stat}[2]{\colorbox{pale}{\parbox{0.9\linewidth}{\color{copper}\Large\bfseries #1\par\color{espresso}\small #2}}\par\vspace{0.5em}}
\newcommand{\qnote}[1]{\par\vspace{0.3em}{\small\color{greyc}#1}}
\newcommand{\tk}[1]{\vfill\noindent\colorbox{copper}{\parbox{\dimexpr\linewidth-2\fboxsep}{\color{white}\bfseries\ Takeaway: #1}}\vspace{0.2em}}
\newcommand{\qframe}[8]{%
\begin{frame}{#1}
\begin{columns}[T]
\column{0.60\textwidth}\centering\includegraphics[width=0.92\linewidth]{#2}
\column{0.37\textwidth}
\stat{#3}{#4}
\stat{#5}{#6}
\qnote{#7}
\end{columns}
\tk{#8}
\end{frame}}

\title{Does a Prague cappuccino get cheaper away from the centre?}
\author{Rishin Kumar}
\institute{Business Analytics, BANL 1100/1101}
\date{}

\begin{document}

% Title ----------------------------------------------------------
{\setbeamercolor{background canvas}{bg=espresso}\setbeamercolor{normal text}{fg=white}\setbeamercolor{itemize body}{fg=white}
\begin{frame}[plain]
\color{white}
\vspace{2.2cm}
{\Huge\bfseries Does a Prague cappuccino get cheaper away from the centre?\par}
\vspace{0.8cm}
{\large\color{copper!60!white} Price, distance and caf\'e type in \Nall{} caf\'es \,\textperiodcentered\, Batch 1\par}
\vspace{1.5cm}
{\small Rishin Kumar \,\textperiodcentered\, Business Analytics, BANL 1100/1101}
\end{frame}}
\note{Good morning, everyone. I'm Rishin. I looked at one question in the class's Prague caf\'e dataset: does a cappuccino get cheaper as you move away from the centre? I planned ten analyses in advance. For each finding, I'll tell you which test backs it up and why I trust it. The data: \Nall{} caf\'es after cleaning, visited by ten students. The average cappuccino costs \PriceMean{} crowns, and the typical caf\'e is \DistMedian{} kilometres from the Orloj. Everything I say is ``in this sample'': it is a convenience sample, not a random sample of Prague. Cleaning is all done in an R script, so anyone can rerun it and get my numbers.}

% One file per question ------------------------------------------
%% A1.tex -- slide for question A1; the graph is figures/A1_figure.png
\qframe{A1: Is the average price different from my prior guess?}{A1_figure.png}
  {\AoneMean{} CZK}{mean cappuccino, 95\% CI \CI{\AoneLo}{\AoneHi}}
  {$d$ = \AoneD}{one-sample $t$-test vs.\ \AoneBench{} CZK, \AoneP}
  {Wilcoxon and BCa bootstrap agree. Prior guess: \AoneBench{} CZK.}
  {A typical cappuccino costs about 83 CZK.}
\note{A1. What does a cappuccino cost? About \AoneMean{} crowns, with a 95 percent interval from \AoneLo{} to \AoneHi. [Say your own prior guess here, when you wrote it down, and the result. The file currently uses \AoneBench{} crowns as a placeholder.] I used a one-sample $t$-test. Defence: the price is only mildly skewed and the sample is large, and the Wilcoxon test and a bootstrap interval give the same answer.}

%% A2.tex -- slide for question A2; the graph is figures/A2_figure.png
\qframe{A2: Is Prague 1 dearer than the rest of Prague?}{A2_figure.png}
  {+\AtwoDiff{} CZK}{Prague 1 $-$ rest, 95\% CI \CI{\AtwoLo}{\AtwoHi}}
  {$d$ = \AtwoD}{Welch $t$-test, \AtwoP{} (one-tailed)}
  {Rank-sum check: \AtwoWP, shift \AtwoHL{} CZK. The tests only partly agree.}
  {Prague 1 is probably a few crowns dearer, but the tests only partly agree.}
\note{A2. Prague 1 averages \AtwoMa{} crowns against \AtwoMb{} elsewhere: \AtwoDiff{} crowns more, interval \AtwoLo{} to \AtwoHi, a small-to-medium effect. Welch's test, one-tailed as planned, gives \AtwoP. But the rank-based Wilcoxon test gives \AtwoWP{} and a typical shift of only \AtwoHL{} crowns. Defence: I chose Welch because the question is about averages and both groups are large, but I report the disagreement rather than hide it. A few very dear Prague 1 caf\'es pull the average up, so the honest summary is: probably 3 to 6 crowns dearer, under ten percent of the price.}

%% A3.tex -- slide for question A3; the graph is figures/A3_figure.png
\qframe{A3: How much more is a cappuccino than an espresso?}{A3_figure.png}
  {+\AthreeM{} CZK}{cappuccino $-$ espresso, 95\% CI \CI{\AthreeLo}{\AthreeHi}}
  {$d_z$ = \AthreeDz}{paired $t$-test, \AthreeP{} (one-tailed)}
  {About \AthreePctEsp\% on top of the espresso price. Wilcoxon agrees.}
  {A cappuccino costs about 28\% more than an espresso.}
\note{A3. In the same caf\'e, the cappuccino costs \AthreeM{} crowns more than the espresso, interval \AthreeLo{} to \AthreeHi: about \AthreePctEsp{} percent more. Paired $t$-test, one-tailed because I expected the cappuccino to be dearer, \AthreeP. Defence: pairing compares each caf\'e with itself, so differences between caf\'es cannot distort it, and the Wilcoxon test, which makes no normality assumption, agrees.}

%% A4.tex -- slide for question A4; the graph is figures/A4_figure.png
\qframe{A4: Do prices differ across the six districts?}{A4_figure.png}
  {$F(\AfourDfa,\,\AfourDfb)$ = \AfourF}{Welch ANOVA, \AfourP}
  {$\eta^2$ = \AfourEta}{95\% CI \CI{\AfourEtaLo}{\AfourEtaHi}}
  {Kruskal--Wallis agrees (\AfourKWP). No pair of districts differs after adjustment.}
  {No clear price differences between districts.}
\note{A4. Welch's ANOVA gives \AfourP{} and the Kruskal--Wallis test \AfourKWP. Districts explain only about 6 percent of the variation. No single pair of districts differs once I correct for comparing many pairs. Defence: with fifteen pairs some will look different by chance, which is why I used Games--Howell and Dunn--Holm corrections. The biggest gap, \AfourBestPair, is \AfourBestDiff{} crowns, but its interval includes zero.}

%% A5.tex -- slide for question A5; the graph is figures/A5_figure.png
\qframe{A5: Is caf\'e type linked to being in Prague 1?}{A5_figure.png}
  {\AfiveFisherP}{Fisher's exact test}
  {$V$ = \AfiveV}{Cram\'er's $V$, 95\% CI \CI{\AfiveVLo}{\AfiveVHi}}
  {\AfiveKInt{} of \AfiveNInt{} international chains are in Prague 1.}
  {International chains cluster in Prague 1; suggestive, not conclusive.}
\note{A5. \AfiveKInt{} of the \AfiveNInt{} international chains, about \AfiveShareInt{} percent, are in Prague 1, compared with about \AfiveShareInd{} percent of independents. Fisher's exact test gives \AfiveFisherP. Defence: I used Fisher's test instead of chi-square because one expected count is small. The pattern is suggestive, not conclusive.}

%% A6.tex -- slide for question A6; the graph is figures/A6_figure.png
\qframe{A6: Does price fall with distance from the Orloj?}{A6_figure.png}
  {$r$ = \AsixR}{Pearson, 95\% CI \CI{\AsixRLo}{\AsixRHi}}
  {\AsixP}{one-tailed; distance shares \AsixRsq\% of price variance}
  {Spearman $\rho$ = \AsixRho{} and Kendall $\tau$ = \AsixTau{} agree.}
  {Farther from the centre, prices are slightly lower.}
\note{A6. The scatterplot shows a downward trend, but a weak one. Pearson's correlation is \AsixR, one-tailed \AsixP. Distance explains under three percent of the variation in price. Defence: I decided the direction in advance, which is why a one-tailed test is fair. Distance is skewed, so I also ran Spearman and Kendall correlations, \AsixRho{} and \AsixTau: the same story.}

%% A7.tex -- slide for question A7; the graph is figures/A7_figure.png
\qframe{A7: By how much does price change per kilometre?}{A7_figure.png}
  {\AsevenPct\% per km}{log-linear, robust 95\% CI \CI{\AsevenPctLo}{\AsevenPctHi}}
  {\AsevenLogP}{two-tailed, HC3 robust standard errors}
  {About \AsevenB{} CZK per km. \AsevenNinfl{} influential caf\'es: without them \AsevenBinfl{} CZK per km.}
  {About 2\% cheaper per km: a small, uncertain effect.}
\note{A7. Each extra kilometre goes with a price change of about \AsevenPct{} percent, roughly \AsevenB{} crowns. The interval runs from \AsevenPctLo{} to \AsevenPctHi{} percent, so it just touches zero. Defence: this test is two-tailed with robust standard errors, so it is stricter than A6, which is why it is less conclusive. Eleven influential caf\'es matter: without them the slope is about \AsevenBinfl{} crowns, still negative. This shows an association only: rent, tourist traffic and caf\'e type all move with distance, so the data cannot say why prices fall.}

%% A8.tex -- slide for question A8; the graph is figures/A8_figure.png
\qframe{A8: Does the distance effect survive controlling for caf\'e type?}{A8_figure.png}
  {\AeightBa{} $\to$ \AeightBb{}}{distance slope (CZK/km) without $\to$ with caf\'e type}
  {$\Delta R^2$ = \AeightDeltaR}{nested $F$-test, \AeightP}
  {Local chains charge +\AeightLocal{} CZK over independents at the same distance.}
  {Distance still matters after caf\'e type, and chains charge more.}
\note{A8. Once I add caf\'e type, the distance slope barely moves, from \AeightBa{} to \AeightBb{} crowns per kilometre. Caf\'e type itself improves the model, nested $F$-test \AeightP, and local chains charge about \AeightLocal{} crowns more than independents at the same distance. Defence: the variance inflation factors are close to one, so distance and type carry separate information.}

%% A9.tex -- slide for question A9; the graph is figures/A9_figure.png
\qframe{A9: Does the distance slope differ for independents and chains?}{A9_figure.png}
  {\AnineInt{} CZK/km}{interaction, robust 95\% CI \CI{\AnineIntLo}{\AnineIntHi}}
  {\AnineP}{nested $F$-test; $\Delta R^2$ = \AnineDeltaR}
  {Slopes: independents \AnineSlopeInd, chains \AnineSlopeCh{} CZK/km. No clear difference.}
  {No clear evidence that the distance effect differs for chains.}
\note{A9. The slope is \AnineSlopeInd{} crowns per kilometre for independents and \AnineSlopeCh{} for chains. That looks different, but the interaction has \AnineP{} and an interval from \AnineIntLo{} to \AnineIntHi. Defence: no clear difference is not the same as evidence the slopes are equal. The data are simply too noisy to tell.}

%% A10.tex -- slide for question A10; the graph is figures/A10_figure.png
\qframe{A10: How does the chance of an above-median price change per km?}{A10_figure.png}
  {\AtenAME{} pp per km}{average marginal effect, 95\% CI \CI{\AtenAMELo}{\AtenAMEHi}}
  {OR = \AtenOR}{logit, \AtenP}
  {Outcome: price above 83 CZK (the brief's median).}
  {Each extra km lowers the chance of an above-median price by about 5 points.}
\note{A10. I turned price into a yes-or-no question: is a cappuccino above 83 crowns, the median in the project brief? In a logistic regression each extra kilometre changes that probability by an average of \AtenAME{} percentage points, interval \AtenAMELo{} to \AtenAMEHi, \AtenP; the odds ratio is \AtenOR. Defence: I report percentage points because they are easier to read than odds. It points the same way as everything else, but the interval includes zero, so I am not overselling it. To close: I ran \Ntests{} tests in total, so one borderline $p$-value is not a finding, which is why I focus on the ten planned analyses and on effect sizes and intervals. This is a convenience sample, the data are observational, and caf\'es in the same district are not fully independent. Thank you; I am happy to take questions.}


\end{document}
